\textbf{结论名称：} 函数值域——分离常数法（一次分式型）

\textbf{适用条件：} 函数形如 $y=\dfrac{ax+b}{cx+d}$，其中 $c\neq0$ 且 $ad\neq bc$（即分子分母不成常数倍），定义域为 $x\neq -\dfrac{d}{c}$。

\textbf{变量说明：} $a,b,c,d\in\mathbb{R}$ 为常数，$c\neq0$，$ad\neq bc$；自变量 $x\in\mathbb{R}\setminus\{-d/c\}$。

\textbf{核心公式：}
\[
y=\frac{ax+b}{cx+d}=\frac{a}{c}+\frac{bc-ad}{c(cx+d)},\qquad \boxed{\{y\in\mathbb{R}\mid y\neq \frac{a}{c}\}}
\]

\textbf{使用场景：} 求一次分式函数的值域、确定图像的水平渐近线、求反函数定义域。

\textbf{简单例子：} 求 $y=\dfrac{2x+3}{x-1}$ 的值域。  
解：$y=\dfrac{2(x-1)+5}{x-1}=2+\dfrac{5}{x-1}$。  
由 $\dfrac{5}{x-1}\neq0$，得 $y\neq2$，故值域为 $\{y\mid y\neq2\}$。